\begin{table}[H]
\centering
\captionsetup{justification=centering}
\caption{Resultados pontuais do modelo multivariado para o total de ocupados - 4\textordmasculine{} trimestre de 2024}
\label{tab:pontual_ocup}

\renewcommand{\arraystretch}{0.8}
\resizebox{\textwidth}{!}{%
\begin{tabular}{@{}p{2.8cm}cccccc@{}}
\toprule
\textbf{Estrato Geográfico} & $\hat{y}$ & $CV(\hat{y})$ & IC 95\% & $\hat{\theta}$ & $CV(\hat{\theta})$ & IC 95\% \\
\midrule
01 - Belo Horizonte & 1394,54 & 1,76\% & [1346,53 ; 1442,54] & 1391,11 & 1,63\% & [1346,66 ; 1435,56] \\
\raisebox{0em}{\parbox[t]{4.3cm}{02 - Entorno e Colar \\[-2pt] Metropolitano de BH}} & 1923,78 & 2,07\% & [1845,68 ; 2001,89] & 1929,28 & 1,89\% & [1857,79 ; 2000,77] \\
03 - Sul de Minas & 1429,66 & 4,45\% & [1304,90 ; 1554,42] & 1433,20 & 3,51\% & [1334,72 ; 1531,68] \\
04 - Triângulo Mineiro & 1582,25 & 4,72\% & [1435,74 ; 1728,75] & 1523,33 & 3,89\% & [1407,31 ; 1639,34] \\
05 - Zona da Mata & 1075,88 & 4,49\% & [981,15 ; 1170,60] & 1043,06 & 3,78\% & [965,68 ; 1120,44] \\
06 - Norte de Minas & 1270,40 & 4,39\% & [1161,19 ; 1379,62] & 1231,96 & 3,82\% & [1139,62 ; 1324,31] \\
07 - Vale do Rio Doce & 1024,62 & 5,35\% & [917,08 ; 1132,16] & 1030,74 & 3,71\% & [955,69 ; 1105,78] \\
08 - Central & 1374,72 & 4,51\% & [1253,32 ; 1496,11] & 1389,45 & 3,66\% & [1289,85 ; 1489,05] \\
\bottomrule
\end{tabular}%
}
\fonte{Elaboração própria, com base nos dados da PNAD Contínua. Nota: \(\hat{y}\) é a estimativa direta e \(\hat{\theta}\) o sinal estimado pelo modelo multivariado (tendência mais sazonalidade). Valores em milhares de pessoas.}
\end{table}
